\chapter{Domänenspezifische Handles zur statischen Partitionierung und konfliktarmen Ausführung einer ECS-basierten C++-Game-Engine}\label{ch:vorschlag2}

\section{Forschungsfrage}

\begin{quote}
    Inwieweit können domänenspezifische Handles als statische Partitionierungsgrenzen von ECS-Systemen dienen, unterschiedliche Laufzeitbereiche typseitig voneinander abgrenzen und dadurch eine laufzeitgünstige sowie konfliktarme parallele Ausführung unterstützen?
\end{quote}

\section{Methodik}

Die Arbeit betrachtet dabei drei Bereiche:

\begin{enumerate}
    \item \textbf{Architekturgrenzen:}
    Es wird betrachtet, welche Architektur- und Laufzeitgrenzen sich durch domänenspezifische Handles im Typsystem abbilden lassen: Welche sind im gegebenen Kontext sinnvoll?

    \item \textbf{Sichtbarkeit von Systemgrenzen im Scheduling:}
    Es wird untersucht, ob Scheduling-Konflikte durch domänenspezifische Handles bereits statisch sichtbar gemacht oder ausgeschlossen werden können.

    \item \textbf{Messung und Vergleich:}
    Dem empirischen Vergleich liegen folgende Varianten zugrunde:
    \begin{enumerate}
        \item unpartitionierte Welt ohne Parallelisierung,
        \item unpartitionierte Welt mit Parallelisierung,
        \item partitionierte Welt mit domänenspezifischen Handles und Parallelisierung.
    \end{enumerate}
\end{enumerate}

Für die Varianten werden folgende Mess- und Bewertungskriterien herangezogen:

\begin{enumerate}
    \item Laufzeitverhalten, insbesondere Frametime und Update-Zeit der betrachteten Systeme,
    \item Skalierung bei wachsender Objektanzahl,
    \item Kompilierzeit,
    \item Code-Komplexität anhand statischer Code-Analyse,
    \item Anzahl und Art der aus dem statischen Modell ableitbaren Scheduling-Konflikte.
\end{enumerate}

\section{Vorläufige Gliederung}

\begin{enumerate}
    \item Einleitung und Zielsetzung
    \item Grundlagen: ECS, datenorientierte Speicherlayouts und Nebenläufigkeit in Game Engines
    \item \textit{ECS Core}-Modell und Abgrenzung zu \texttt{helios::ecs}
    \item Domänenspezifische Handles als statische Partitionierungsgrenzen
    \item Typisierte Systemsignaturen: Lese-/Schreibmengen
    \item Relationales ECS-Modell als Grundlage komponentenbezogener Lese-/Schreibzugriffe
    \item Konfliktmodell für Systeme und Scheduling
    \item Referenzszenario: Kollisions- und Partikelsysteme
    \item Implementierung der Varianten
    \item Evaluation: Laufzeit, Skalierung, Kompilierzeit und Code-Komplexität
    \item Diskussion
    \item Zusammenfassung und Fazit
\end{enumerate}